\subsection{The spectrum on the endpoint-one surface}

Principal interlacing separates the five positive quotient roots into
two below $a+b$ and three above $b+c$. On endpoint one, this gap
isolates the smallest residual root from the other three.

\begin{theorem}
\label{thm:endpoint-one-spectrum}
Let $2\leq a\leq b\leq c$ be real. The ordered eigenvalues of the
six-support complement quotient, counted with multiplicity, satisfy
\[
 0=\lambda_1<\lambda_2\leq\lambda_3<a+b\leq b+c
 <\lambda_4\leq\lambda_5\leq\lambda_6.
\]
If also $a\geq4$ and $h_C(1)=0$, then
\[
 b+c\geq2a^2-2a+1,
\]
with equality exactly when $b=a$ and $c=(a-1)(2a-1)$. Root one is
simple and is the smallest positive quotient root; all four other
quotient roots are strictly greater than $a$. Their ordered roots satisfy
\[
 a<\mu_1<\min\{c,a+b\},\qquad b+c<\mu_2\leq\mu_3\leq\mu_4.
\]
\end{theorem}

\begin{proof}
Let $S=W^{1/2}CW^{-1/2}$, with $W=\operatorname{diag}(a,b,c,ab,ac,bc)$.
It is symmetric positive semidefinite with a one-dimensional zero kernel.
The principal block on supports $3,12,13,23$ is
\[
 \begin{pmatrix}ab+a+b&-\sqrt{abc}\\-\sqrt{abc}&c\end{pmatrix}
 \mathbin{\oplus}(b)\mathbin{\oplus}(a).
\]
The two-dimensional block has polynomial
$Q(x)=x^2-(ab+a+b+c)x+c(a+b)$, with $Q(a+b)=-ab(a+b)<0$.
Exactly three eigenvalues of this four-dimensional block lie below $a+b$;
interlacing gives $\lambda_3<a+b$.

The principal block on singleton supports $1,2$, shifted by $b+c$, is
\[
 \begin{pmatrix}bc&-\sqrt{ab}\\-\sqrt{ab}&a(c+1)-b\end{pmatrix}.
\]
Its first leading minor is $bc>0$ and its determinant is
$b[a(c^2+c-1)-bc]>0$, since
$a(c^2+c-1)-bc\geq c^2+2c-2>0$.
Both principal eigenvalues therefore exceed $b+c$, so interlacing gives
$\lambda_5,\lambda_6>b+c$. Direct determinant expansion gives
\[
 h_C(b+c)=-abc(b+c+1)
 \bigl[bc(b+c)(a-1)+a(bc-b-c+a)\bigr]<0,
\]
where $bc-b-c+a=(b-1)(c-1)+a-1>0$.
No eigenvalue equals $b+c$. At least three eigenvalues lie below it
and at most four. The negative determinant $(b+c)h_C(b+c)$ requires
an odd number below it, counting multiplicity; hence exactly three
lie below it and $\lambda_4>b+c$.

Now suppose $a\geq4$ and $h_C(1)=0$.
Put $q=b+c$, $z=bc$ and $K=4a^2-2a+1$. The generic determinant is
\begin{align*}
 h_C(1)&=(q-K)z^2-a(a-1)(q+2a-1)z+P_a(q),\\
 P_a(q)&=(a^2-1)q^3+(a^3-a^2-3a+3)q^2
          -3(a-1)^2q-(a-1)^3.
\end{align*}
Since $(b-a)(c-a)\geq0$, we have $z\geq z_0=a(q-a)$. When
$q\leq Q:=2a^2-2a+1$, the coefficient $q-K\leq-2a^2$ is negative,
and the difference between the displayed polynomial at $z$ and at
$z_0$ is
\[
 (z-z_0)\bigl[(q-K)(z+z_0)-a(a-1)(q+2a-1)\bigr]\leq0,
\]
strictly unless $z=z_0$. Set $t=q-a\geq a$. At $z_0$ the value factors as
\[
 \bigl[(a-1)(2a-1)-t\bigr]B(a,t),
\]
where $B=a^3-2a^2t^2+a^2t+a^2+3at-3a+t^2-2t+1$. For
$a=4+m$, $t=a+u$, with $m,u\geq0$, the complete identity
\begin{align*}
 -B(4+m,4+m+u)={}&2m^4+4m^3u+30m^3+2m^2u^2+47m^2u+163m^2\\
 &+16mu^2+179mu+381m+31u^2+222u+323
\end{align*}
proves $B<0$. Thus $q<Q$ gives $h_C(1)<0$. At $q=Q$, equality
requires $z=z_0$, hence $b=a$ by ordering, and $c=(a-1)(2a-1)$.
The factorization also verifies this equality case.

The pair-support principal block is $\operatorname{diag}(c,b,a)$.
The sum bound implies
$b+c-(a^2+2a)\geq(a-4)^2+4(a-4)+1>0$, so
\[
 h_C(a)=-a^2b^2c^2(a^2+2a-b-c)>0.
\]
No eigenvalue equals $a$. The positive determinant
$\det(aI-S)=a h_C(a)$ forces an even number of eigenvalues below
$a$, counted with multiplicity. There are at least two, zero and one,
and at most three by interlacing. There are therefore exactly two:
zero and a simple one. All four remaining roots exceed $a$.

The pair-support block gives $\lambda_3\leq c$. The identity
$h_C(c)=-a^2b^2c^2(c^2+2c-a-b)<0$ excludes equality.
Together with the global pair-sum gap, this proves the residual bounds.
\end{proof}

\begin{corollary}
\label{cor:endpoint-one-divisor-window}
For integer exponents $4\leq a\leq b\leq c$ on $h_C(1)=0$, define
$F(x)=h_C(x)/(x-1)$ and $s=a+b+c$, $p=ab+ac+bc$, $r=abc$.
Then $F$ is a monic integer quartic with $F(0)=rs(p+s)>0$.
Its ordered real roots satisfy
\[
 a<\mu_1<\min\{c,a+b\},\qquad b+c<\mu_2\leq\mu_3\leq\mu_4.
\]
An integer spectrum therefore requires a positive divisor $d$ of $rs(p+s)$ with
\[
 a<d<\min\{c,a+b\},\qquad F(d)=0.
\]
The other three integer roots would each be at least $b+c+1$.
\end{corollary}
\begin{proof}
The determinant gives $h_C(0)=-rs(p+s)$, so division by $x-1$
retains the stated nonzero constant. Apply the rational-root theorem
to the smallest quartic root and apply Theorem~\ref{thm:endpoint-one-spectrum}
to all four roots.
\end{proof}

All four quartic roots exceed $a$; the upper window in the corollary
applies to their smallest root. For example, $(8,8,105)$ has the
quartic factor
\[
 (x-968)(x^3-1138x^2+177145x-1566600),
\]
so another root can exceed both $c$ and $a+b$. Testing the divisor
window is finite for each fixed triple. A surviving candidate still
requires checking the other factors, and the exponent domain is unbounded.

\begin{remark}
The generic residual $[h_C(x)-xh_C(1)]/(x-1)$ is positive on
$[0,4]$ for all real $a,b,c\geq8$, including off the endpoint-one
surface. A separate complete $350$-term positive identity after
$a=8+m$, $b=a+t$, $c=b+u$, $x=4v/(1+v)$ verifies that statement,
including $x=4$. The shorter proof above suffices for the endpoint-one
spectral theorem; the generic residual calculation is retained in the
\href{https://github.com/the-omega-institute/ideal-intersection-laplacian/blob/95d8d9a90d4ef0aa06be8ba5ce70d3355838385c/notes/endpoint-one-spectral-gap.md}{verification note}.
\end{remark}

\begin{remark}[The equality curve]
\label{rem:endpoint-one-equality-curve}
On $c=b(b+2)-a$, put $t=b(b+2)$ and $q=a(t-a)=ac$. Then
$h_C(1)=-(b-1)G(a,b)$, where the plane septic is
\[
 G=(3b-1)q^2+b(b^2+4b-1)q
   -(b^2+2b-1)(b^2+3b-1)(b^3+3b^2+3b-1).
\]
Its smooth projective model is geometrically irreducible of genus ten.
Indeed, the involution $a\mapsto t-a$ has quotient $y^2=D(b)$, where
$D$ is the squarefree degree-eight discriminant of the displayed quadratic
in $q$, so the quotient has genus three. Recovering $a$ adjoins
$\sqrt{t^2-4q}$. The norm of $t^2-4q$ has nine distinct simple zeros
away from the quotient branch points and one simple pole at $b=1/3$.
This nonsquare norm makes the cover connected. Above these points the
cover has ten simple branches; at the two quotient points at infinity,
$t^2-4q$ has even pole order four. Riemann--Hurwitz gives
$2g-2=2(2\cdot3-2)+10=18$. The explicit norm and squarefreeness
identities are retained in the
\href{https://github.com/the-omega-institute/ideal-intersection-laplacian/blob/8ed4f9014f6142949e6af40fd16e18ca304f629a/notes/endpoint-one-equality-curve.md}{curve verification note}.
By Siegel's theorem \cite{hindrysilverman,serre}, this affine equality curve
has finitely many integer points. Their effective determination and
the classification of admissible exponent triples remain open.
\end{remark}

\begin{proposition}
\label{prop:endpoint-one-equality-spectrum}
For real $4\leq a\leq b$ and $c=b(b+2)-a$, the ordered eigenvalues
of the six-support complement quotient satisfy
\[
 0=\lambda_1<\lambda_2<\lambda_3=b<a+b\leq b+c
 <\lambda_4\leq\lambda_5\leq\lambda_6,
\]
and $b$ is simple. If also $h_C(1)=0$, then $\lambda_2=1$ and
$F(x)=(x-b)P_3(x)$, where the monic cubic $P_3$ has all three roots
strictly greater than $b+c$.
\end{proposition}

\begin{proof}
Use the symmetric similarity $S$ from Theorem~\ref{thm:endpoint-one-spectrum}.
Every proper principal block is positive definite: $S$ is positive
semidefinite with a one-dimensional zero kernel whose vector has every
coordinate nonzero. Delete singleton support $2$. The five-support block
$B$ is the direct sum of the scalar $b$ on support $13$ and the block $T$
on $1,3,12,23$. Its characteristic polynomial is
\begin{align*}
 \det(xI-T)&=(x-b)K(x),\\
 K(x)&=x^3-b(b^2+4b+5)x^2\\
 &\quad+\bigl[ac(b^2+1)+b^2(b+2)(b^2+3b+3)\bigr]x-abc(b+3).
\end{align*}
The three roots of $K$ are positive and real. Moreover,
\[
 K(b)=b(b+1)\bigl[a(b-a)(b-2)+ab(b-2)(b+1)+b^2(b+1)(b+2)\bigr]>0.
\]
The pair-support block on $12,23$ is $\operatorname{diag}(c,a)$, so
the largest eigenvalue of $T$ is at least $c>b$. A positive monic cubic
value at $b$ requires one or three roots below $b$; the largest root
rules out three. Thus $B$ has $\beta_2=\beta_3=b$, and interlacing gives
$\beta_2\leq\lambda_3\leq\beta_3$, hence $\lambda_3=b$.

The kernel of $B-bI$ has dimension two. Its support-$13$ unit vector
couples to the deleted singleton by $-\sqrt{abc}\ne0$. Projecting the
full eigenvector equation onto that vector forces the deleted coordinate
to vanish. The remaining eigenvectors lie in $\ker(B-bI)$ and satisfy
one nonzero coupling condition. Their space is one-dimensional, proving
simplicity and $\lambda_2<b<\lambda_4$.

The global gap in Theorem~\ref{thm:endpoint-one-spectrum} gives
$\lambda_4>b+c$. On endpoint one it also gives $\lambda_2=1$, and division
by the simple roots $1,b$ yields the stated cubic.
\end{proof}

This proposition narrows the equality splitting problem: for integer
endpoint-one equality points, all three integer roots of $P_3$ would have
to exceed $b+c$. The
\href{https://github.com/the-omega-institute/ideal-intersection-laplacian/blob/fb1d7c0134ad385931e3b4e43e445695bc4bf566/notes/global-pair-sum-separation.md}{global spectral verification note}
records the determinant identities and exact controls.
